\clearpage
\SourceLeaf{415}%
\begin{center}\large NEUE ALLGEMEINE UNTERSUCHUNGEN\\[4pt] ÜBER\\[4pt] DIE KRUMMEN FLÄCHEN.\\[7pt] [1825.]\end{center}
\noindent\hspace*{1.2em}Zwar besteht der eigentliche Zweck dieses Werks in der Ausarbeitung neuer Lehrsätze über seinen Gegenstand. Doch werden wir anfangs auch das bereits Bekannte entwickeln, teils des Zusammenhangs und der Vollständigkeit wegen, teils weil sich unsere Art der Behandlung von der bisherigen unterscheidet. Aus denselben Gründen werden wir sogar zuerst einige Bemerkungen über Kurven in der Ebene vorausschicken.

\Needspace{5\baselineskip}
\begin{center}1.\end{center}
\noindent\hspace*{1.2em}Um die verschiedenen Richtungen gerader Linien in einer Ebene bequem miteinander vergleichen zu können, denke man sich in derselben Ebene um einen beliebigen Mittelpunkt einen Kreis mit dem Radius $1$ beschrieben. Die Lage des parallel zu einer vorgegebenen Linie gezogenen Radius dieses Kreises veranschaulicht dann die Lage jener Linie. Die Neigung zweier Geraden zueinander wird durch den Winkel der beiden darstellenden Radien oder durch den Bogen zwischen ihren Endpunkten gemessen. Wo eine genaue Bestimmung nötig ist, muss selbstverständlich bei jeder\space%
